% The comparator and the obstruction --- plan \S4.2.
% Sources: draft \S2 (02_why_constant_release_fails.tex), rewritten from the
% original Ch.4 \S9 attempt.  The canonical parameter trio is introduced
% here and used by name for the rest of the chapter.

\section{The comparator and the obstruction}
\label{p:sec:comparator}

The single-cell theory is complete; what it implies for a population of
infected cells is not. The obstruction shows itself most clearly against the
standard model this chapter is meant to improve upon, which is where we
begin.

\subsection{The budding cell}
\label{p:sec:classical}

In the simplified, target-cell-constant Basic Model of Viral Replication
\cite{mclean1993balance,nowak1996population,perelson1996hiv}, infected
cells $\Icell$ and free particles $\Vfree$ obey
\begin{equation}
\ddt{\Icell} = \gamma T\,\Vfree - d_{\Icell}\,\Icell,
\qquad
\ddt{\Vfree} = p\,\Icell - c\,\Vfree,
\label{p:eq:bmvr-classical}
\end{equation}
with target cells $T$ fixed, infection rate $\gamma$, virion clearance $c$,
and the two constants that concern us: $p$, the per-cell release rate, and
$d_{\Icell}$, the removal rate. Both are phenomenological, parameters to be
fitted rather than quantities derived from anything inside the cell. The
basic reproduction number is
\begin{equation}
R_0=\frac{\gamma T\,p}{c\,d_{\Icell}}.
\label{p:eq:R0classical}
\end{equation}
Every cell, regardless of its age or its contents, releases at rate $p$ and
dies at rate $d_{\Icell}$.

That is not an arbitrary caricature. It is one kind of cell described
exactly, and it is worth setting out in full, because it is the comparator
the whole chapter is built against. A budding infected cell experiences two \emph{independent} event
types: release of one particle, as a Poisson process of rate $p$, and death,
as a Poisson process of rate $d_{\Icell}$; the two run side by side until
the cell dies. Its survival and cumulative release are
\begin{equation}
I_{\rm bud}(\alpha)=\ee^{-d_{\Icell}\alpha},\qquad
V_{\rm bud}(\alpha)=\frac{p}{d_{\Icell}}\lrb{1-\ee^{-d_{\Icell}\alpha}},
\qquad
\dda{V_{\rm bud}}{I_{\rm bud}}=-\frac{p}{d_{\Icell}}:
\label{p:eq:bud-cell}
\end{equation}
the mean lifetime is $1/d_{\Icell}$, the lifetime release $p/d_{\Icell}$,
and the total released is geometric on $\{0,1,2,\ldots\}$ with
$\pr{\Kb_{\rm bud}=k}=\frac{d_{\Icell}}{d_{\Icell}+p}
\bigl(\frac{p}{d_{\Icell}+p}\bigr)^k$ --- the successes before the
terminating death in a sequence of events each of which is a release with
probability $p/(d_{\Icell}+p)$. All of it sits in the pair $(p,d_{\Icell})$. That is why
\cref{p:eq:bmvr-classical} is a faithful one-cell limit for a budding
pathogen, and it is this cell that \cref{p:sec:flooding} matches against a
bursting one.

The structural contrast with bursting is the engine of everything that
follows. In budding, release and death are independent: a cell may release
many particles and then die with its contents lost, or release nothing at
all. In bursting they are the \emph{same} event: whatever has accumulated is
released at once, and the cell's productive life ends in that instant. The
detailed single-cell comparison is in \ChDist{}; here the point is that
\cref{p:eq:bmvr-classical} encodes the budding picture exactly, and the
question is what replaces it when release comes in terminal batches.

\subsection{The synchronous cohort}
\label{p:sec:cohort}

Take a clean special case: $N_0$ cells all infected at the same instant,
each running an independent copy of the single-cell process, with no
reabsorption or extracellular replication. The number still productively
infected at time $t$ is $N_0\,\Ifix(t)$, exactly in expectation, and the
expected cumulative extracellular count is
\begin{equation}
\Vfree(t)=N_0\,V(t),
\label{p:eq:cohort}
\end{equation}
since $V(t)$ is the expected release from one cell. The release \emph{rate}
is then $N_0V'(t)=N_0\delta K(t)$: it depends on the age of the cohort, and
it is not constant. The classical model's $p\Icell$ can mimic such a cohort
only by accident, and only briefly (\cref{p:fig:constant-fails}).

\begin{figure}[tbp]
\centering
\includegraphics[width=0.94\textwidth]{N4b_1_constant_release_fails.pdf}
\caption{A synchronous cohort releases on a schedule, not at a rate. Both
schedules have lifetime yield $V_\infty=11$ and mean productive lifetime
$\mean{T_{\rm prod}}=2.398$, with $d_{\Icell}=1/\mean{T_{\rm prod}}$,
$p=V_\infty d_{\Icell}$, $(\lambda,\mu,\delta)=(1,0,0.1)$ and $c=1$: as
fairly matched as two constants allow. (a) Matched yield, different
release timing: the release kernel $g(\alpha)=\delta K(\alpha)$ still lags
the budding cell's $p\ee^{-d_{\Icell}\alpha}$, because a bursting cell must
first accumulate load. (b) Linear clearance preserves that lag, carrying it
into the free-particle response.}
\label{p:fig:constant-fails}
\end{figure}

\subsection{The proposal that cannot work}
\label{p:sec:obstruction}

The natural thought is to replace the constant $p$ in
\cref{p:eq:bmvr-classical} by a constant built from the intracellular
theory: the quasi-stationary mean,
\begin{equation}
p\ \stackrel{?}{=}\ \langle X\rangle_{\rm QS}=\frac{a}{a-1}.
\label{p:eq:naive-p}
\end{equation}
The intuition is that a long-lived cell hovers at the quasi-stationary
level, so its bursting should look like release at rate
$\langle X\rangle_{\rm QS}$. The obstruction is dimensional: $p$ has units
of particles per unit time and $\langle X\rangle_{\rm QS}$ is a pure count,
so \cref{p:eq:naive-p} introduces a spurious factor of time. Worse, when
$\mu=0$ one has $\langle X\rangle_{\rm QS}=V_\infty$ exactly, so the
proposal substitutes a cell's \emph{total lifetime output} for its
\emph{output rate}; \cref{p:fig:overlay-naive} shows what that costs.

Two repairs suggest themselves, and both fail for related reasons. First,
average the burst size over the burst-time density $\delta J$ of \ChDist{}:
\begin{equation}
\int_0^\infty\delta J(t)\,\mean{\Kb\mid\tau=t}\,\dd t
=\delta\int_0^\infty K(t)\,\dd t
=V_\infty,
\label{p:eq:naive-repair}
\end{equation}
using \cref{p:eq:sizebias} and $V'=\delta K$ --- again a count, not a rate.
Second, weight the load-dependent flux by the burst-time density,
$\delta\int\delta J(t)\,\mean{X_t^2\mid\cdot}\,\dd t$: the dimensions are
right, but the measure is wrong. The burst-time density describes \emph{when
a single cell bursts}; a population release rate needs an average over the
\emph{ages of the cells currently infected}, and in an asynchronous
infection those are not the same. No reweighting of a single-cell count by a
single-cell density can fix the substitution, because the difficulty is not
in the single-cell theory at all. It is in the population bookkeeping.

What survives the failed proposal is two facts, not a substitute constant.
The expected release flux of a cell of age $\alpha$ is $\delta K(\alpha)$,
and at any host time the infected population is a mixture of ages, one
cohort per past moment of infection. The population model must therefore
track cell age. That is the renewal construction of
\cref{p:sec:renewal}.

\subsection{Three regimes, used throughout}
\label{p:sec:trio}

Three intracellular parameter triples recur throughout, chosen to straddle
the flooding boundary $L=1$ of \cref{p:sec:flooding}: one on each side, one
exactly on it. They are named here and used by name thereafter, and
\cref{p:tab:trio} collects what follows from them by
\crefrange{p:eq:roots}{p:eq:lifetime}.

\begin{table}[htbp]
\centering
\renewcommand{\arraystretch}{1.12}
\small
\begin{tabular}{@{}lccccccc@{}}
\toprule
$(\lambda,\mu,\delta)$ & $a$ & $b$ & $L$ & $V_\infty$
 & $\mean{T_{\rm prod}}$ & $d_{\Icell,{\rm eff}}(0)$ & $p_{\rm eff}(0)$ \\
\midrule
$(1,0,0.1)$    & $1.100$ & $0$     & $1.10$ & $11.00$ & $2.398$ & $0.417$ & $4.59$ \\
$(1,0.5,1/3)$  & $1.500$ & $1/3$   & $1.00$ & $2.00$  & $1.099$ & $0.910$ & $1.82$ \\
$(1,0.9,0.1)$  & $1.316$ & $0.684$ & $0.42$ & $1.32$  & $1.426$ & $0.701$ & $0.92$ \\
\bottomrule
\end{tabular}
\caption{The three regimes used throughout, and what follows from them: the
first has $L>1$, the second sits on the boundary $L=1$, the third has $L<1$.
The last two columns are the $r=0$ endpoints of the effective-parameter map
of \cref{p:sec:projection}. The verification suite of
\cref{p:app:verification} uses these three with $(1,0.2,0.05)$ and
$(1,0,0.5)$; a wider spread suits testing than exposition.}
\label{p:tab:trio}
\end{table}

\FloatBarrier
